
فرض کنید دو تابع یک‌به‌یک $f: A \to B$ و $g: B \to A$ داریم. می‌خواهیم ثابت کنیم یک تابع یک‌به‌یک و پوشا از $A$ به $B$ وجود دارد.
برای اثبات، سعی می‌کنیم با استفاده از توابع $f$ و $g$، یک تابع جدید مثل $h: A \to B$ بسازیم که هم یک‌به‌یک باشد و هم پوشا.
برای هر عضو $x \in A$، می‌توانیم مسیر حرکت آن را به سمت عقب دنبال کنیم تا ببینیم این عضو از کجا آمده است. زنجیره عقب‌گرد به این صورت است:
\[
\dots \xrightarrow{g} x_2 \xrightarrow{f} y_1 \xrightarrow{g} x_1 \xrightarrow{f} y_0 \xrightarrow{g} x
\]

اگر این زنجیره را برای هر عضو تا جایی که ممکن است به عقب ادامه دهیم، سه حالت پیش می‌آید:
\begin{enumerate}
    \item زنجیره در مجموعه $A$ متوقف می‌شود (یعنی به عضوی می‌رسیم که تصویر هیچ عضوی از $B$ تحت تابع $g$ نیست). مجموعه این اعضا را $A_A$ می‌نامیم.
    \item زنجیره در مجموعه $B$ متوقف می‌شود (یعنی به عضوی می‌رسیم که تصویر هیچ عضوی از $A$ تحت تابع $f$ نیست). مجموعه این اعضا را $A_B$ می‌نامیم.
    \item زنجیره هیچ‌وقت تمام نمی‌شود و تا بی‌نهایت به عقب می‌رود یا دور می‌زند. مجموعه این اعضا را $A_\infty$ می‌نامیم.
\end{enumerate}

واضح است که هر عضو از $A$ دقیقاً در یکی از این سه گروه قرار می‌گیرد. پس مجموعه $A$ به سه بخش جدا از هم افراز می‌شود:
\[
A = A_A \cup A_B \cup A_\infty
\]
به همین ترتیب، مجموعه $B$ را هم می‌توان به سه بخش متناظر افراز کرد:
\[
B = B_A \cup B_B \cup B_\infty
\]

حالا با توجه به یک‌به‌یک بودن توابع اولیه، رفتار آن‌ها روی این بخش‌ها به شکل زیر است:
\begin{itemize}
    \item تابع $f$ مجموعه $A_A$ را دقیقاً و به طور یک‌به‌یک و پوشا به $B_A$ می‌نگارد.
    \item تابع $f$ مجموعه $A_\infty$ را دقیقاً و به طور یک‌به‌یک و پوشا به $B_\infty$ می‌نگارد.
    \item تابع $g$ مجموعه $B_B$ را به $A_B$ می‌نگارد. چون $g$ یک‌به‌یک است، معکوس آن یعنی $g^{-1}$ وجود دارد و مجموعه $A_B$ را به طور یک‌به‌یک و پوشا به $B_B$ می‌برد.
\end{itemize}

حالا تابع نهایی $h: A \to B$ را به صورت تکه‌ای روی این بخش‌ها تعریف می‌کنیم:
\[
h(x) = \begin{cases} 
f(x) & \text{if } \ \ x \in A_A \cup A_\infty \\ 
g^{-1}(x) & \text{if } \ \ x \in A_B 
\end{cases}
\]

در نهایت نشان می‌دهیم که $h$ یک‌به‌یک و پوشا است:

\begin{enumerate}
    \item \textbf{یک‌به‌یک بودن:} 
    فرض کنید $h(x_1) = h(x_2)$. 
    \begin{itemize}
        \item اگر $x_1$ و $x_2$ هر دو در $A_A \cup A_\infty$ باشند، چون $f$ یک‌به‌یک است، نتیجه می‌شود $x_1 = x_2$.
        \item اگر هر دو در $A_B$ باشند، چون $g^{-1}$ یک‌به‌یک است، نتیجه می‌شود $x_1 = x_2$.
        \item اگر یکی در $A_A \cup A_\infty$ و دیگری در $A_B$ باشد، تصویر اولی در $B_A \cup B_\infty$ و تصویر دومی در $B_B$ می‌افتد. چون این مجموعه‌ها هیچ اشتراکی ندارند، غیرممکن است تصاویر آن‌ها برابر شود.
    \end{itemize}
    پس $h$ یک‌به‌یک است.

    \item \textbf{پوشا بودن:} 
    باید بررسی کنیم آیا برد $h$ کل مجموعه $B$ را پر می‌کند یا خیر:
    \[
    h(A) = h(A_A \cup A_\infty) \cup h(A_B) = f(A_A \cup A_\infty) \cup g^{-1}(A_B)
    \]
    با توجه به روابط بخش‌های متناظر داریم:
    \[
    h(A) = (B_A \cup B_\infty) \cup B_B = B
    \]
    پس تمام اعضای $B$ پوشش داده می‌شوند و تابع پوشا است.
\end{enumerate}

